\documentclass[11pt]{article}
\usepackage[a4paper,margin=1in]{geometry}
\usepackage[utf8]{inputenc}
\usepackage[T1]{fontenc}
\usepackage{amsmath,amssymb}
\usepackage{hyperref}
\usepackage{booktabs}

\title{\textbf{Crossover Helps Only Where the Search Space Has Building Blocks:\\
A Trap-Function Control and a Real NAS Negative}}
\author{Vasili Gavrilov\thanks{Independent researcher. Physics and computer
science by education; a professional software engineer with a background in
numerical optimization and machine learning, including a genetic-algorithm
optimizer built for industrial use in 2002. Reference implementation:
\url{https://github.com/Anode1/SMBPANN}.}}
\date{2026}

\begin{document}
\maketitle

\begin{abstract}
Evolutionary neural architecture search recombines architectures with crossover, but random
search is a strong baseline and it is unclear whether crossover adds anything of its own. We
test one rule: crossover helps when the search space is built from separable blocks that
recombination can assemble. As a positive control, on deceptive trap functions (independent
blocks by design) crossover beats mutation, and its advantage grows with the number of blocks,
from 1.8 times at four blocks to sixteen blocks, where mutation reaches the optimum $28$ times
in $200$ runs and crossover $199$. On NAS-Bench-101 ($423{,}624$ tabulated cells, so fitness is
a table lookup and thousands of paired runs are cheap), no method meaningfully beats random:
random and all four crossovers land within a few tenths of a percent of each other, inside the
benchmark's own training noise. Thousands of paired seeds still resolve a consistent ordering,
and it does not favor structure. Every crossover that recombines the cell's structure (wiring
versus labeling, node roles, paths) is consistently worse than a plain structure-blind uniform
crossover, and two fall slightly below random at large budget. So the cell space has no
building-block structure to exploit; the small, consistent edge of uniform crossover over
random comes from diversity, not structure, which we confirm by measuring population diversity
and by a population-size sweep. The result is the familiar one, that random search is hard to
beat, now with a mechanism and a positive control that shows when recombination does help.
Every number regenerates from one command.
\end{abstract}

\section{Introduction}

Neural architecture search (NAS) automates a network's shape; methods differ mainly in search
strategy: reinforcement learning \cite{zoph2017}, gradient relaxation \cite{darts2019},
Bayesian optimization, random search, and evolution \cite{elsken2019}. Evolutionary NAS mutates
and recombines architectures, training each candidate by back-propagation \cite{real2019}. Two
facts motivate us: plain random search is a strong baseline \cite{bergstra2012,li2019}, and the
leading evolutionary method drops crossover entirely, keeping only mutation with aging
\cite{real2019}.

Crossover \cite{holland1975} is meant to recombine \emph{building blocks}, partial solutions
that can be assembled across parents. So the real question is not whether crossover helps, but
whether the space has building blocks to recombine. We answer with one rule and two testbeds.
A positive control fixes what the rule means: deceptive trap functions, which are separable
into blocks by construction, where crossover must and does win. Then the real test:
NAS-Bench-101, whose cell has two natural dimensions, a wiring dimension and a labeling
dimension, plus a path view. We ask whether a crossover that respects that structure beats one
that ignores it. As a hypothesis, $H_1$ says a crossover that recombines the cell's building
blocks (its wiring and labeling, its node roles, its paths) beats a structure-blind uniform
crossover; the null $H_0$ says it does not, because the real cell space carries no such
recombinable structure. The intuition behind $H_1$ is strong, and our trap control shows the
effect it predicts is real when blocks exist. Yet on NAS-Bench-101 the data does not merely
fail to reject $H_0$: the structure-aware operators are consistently worse than the
structure-blind one, and at large budget worse than random search. The recombination logic
comes from a self-modifying back-propagation engine that grew out of the author's 1997 seminar
thesis \cite{thesis1997}; on NAS-Bench-101 we move fitness onto the table so training noise
cannot confound the comparison.

\section{Positive control: crossover on separable building blocks}

Before the real space, we fix what a building-block advantage looks like on a problem
\emph{designed} to have one. The genome is $M$ concatenated deceptive trap blocks of $K=4$
bits: a block scores $K$ if all its bits are $1$, else $K-1$ minus its number of ones. That
gradient lures a bit-flipping search toward all-zero, so the all-ones optimum is reachable only
by assembling correctly-solved blocks, which is what one-point crossover recombines and a
mutation-only search cannot climb to \cite{deb1993}. We run elitism with tournament selection
and one-point crossover ($p=0.9$) against the identical search with mutation only (population
$1000$, $200$ seeds, cap $500$ generations), and scan $M$.

\begin{table}[t]
\centering
\caption{Concatenated deceptive traps ($K=4$ bits per block), mean generations to the global
optimum and how many of $200$ seeds reach it. Crossover's advantage over mutation grows with
the number of blocks $M$, the signature of a genuine building-block effect, until mutation
essentially fails and crossover still solves it.}
\label{tab:bb}
\begin{tabular}{lccc}
\toprule
blocks $M$ & mutation-only & crossover $+$ mutation & speedup \\
\midrule
$4$  & $9.2$ gens, $200/200$  & $5.1$ gens, $200/200$  & $1.8\times$ \\
$8$  & $70.8$ gens, $200/200$ & $14.2$ gens, $200/200$ & $5\times$ \\
$12$ & $257$ gens, $185/200$  & $22.6$ gens, $200/200$ & $11\times$ \\
$16$ & $363$ gens, \textbf{$28/200$} & $42.3$ gens, \textbf{$199/200$} & mutation fails \\
\bottomrule
\end{tabular}
\end{table}

Table~\ref{tab:bb} is the control: when the space is separable into building blocks, crossover
is not marginally but decisively better, and the gap widens with the amount of structure until
mutation alone cannot solve the problem at all. The operator works. The question for a real
search space is whether it supplies such blocks.

\section{The real space and its fitness oracle}

A NAS-Bench-101 cell is a directed acyclic graph on up to seven nodes: node~0 is the input,
node~$n{-}1$ the output, and the interior nodes are labeled
\texttt{conv1x1}, \texttt{conv3x3}, or \texttt{maxpool3x3}. It has two orthogonal dimensions,
its \emph{wiring} (the adjacency) and its \emph{labeling} (the ops), and a derived path view.
The benchmark \cite{nasbench101} tabulates the trained CIFAR-10 accuracy of all $423{,}624$
cells, so fitness is a lookup and no network is trained during search.

The benchmark removes duplicate \emph{isomorphic} graphs, so a recombined child is usually a
relabeling of a catalogued cell rather than a verbatim match, and lookup needs a canonical
form. We compute it by brute force: prune the nodes that lie on no input-to-output path, then
take the lexicographically minimum adjacency-plus-labeling encoding over all $\le5!=120$
permutations of the interior nodes. This is exact and needs no library. A self-test confirms
it: over $5{,}000$ cells both the verbatim graph and a random relabeling map to the same
accuracy ($5000/5000$ each), and the $423{,}624$ rows collapse to exactly $423{,}624$ canonical
keys, with no collisions. The same canonicalization gives the alignment operator its node
correspondence. The table is extracted from the official \texttt{nasbench\_only108.tfrecord} by
a small C decoder; best test $94.32\%$ and best validation $95.06\%$ reproduce the published
optima.

\section{Operators}

Every individual is a seven-node upper-triangular cell. All arms share one population loop, one
light mutation (flip an edge or re-roll an op), the same elitism, and the same query budget;
only the crossover differs. Selection is on validation accuracy, and we report test accuracy.
Four crossovers run against a random-search control:

\begin{itemize}
\item \textbf{flat}: per-cell uniform crossover. Each adjacency bit and each op comes from a
uniformly chosen parent. It is structure-blind and maximally diverse. It also ignores that node
$k$ need not play the same role in both parents (the competing-conventions problem).
\item \textbf{axis}: dimension-factored. The whole wiring comes from one parent, the whole
labeling from the other.
\item \textbf{aligned}: role-aligned. Match the second parent's interior nodes to the first by
role (op and in/out degree), then recombine node modules, each node's incoming edges and its op
taken together from one aligned parent. This removes the competing-conventions problem
\cite{neat2002,modxover2021}.
\item \textbf{path}: subgraph transplant. Recombine whole input-to-output paths, each carrying
its own parent's ops, added greedily under the nine-edge cap. This works along the graph
dimension.
\end{itemize}

\textbf{Matched compute, paired.} Each arm gets a fixed query budget; the random control spends
the same budget on independent random cells; a child that prunes to an invalid or over-budget
graph is a spent query scored zero. In each trial the control and all arms share the same seed,
so they start from an identical population and a trial is a paired observation. We report paired
win/loss counts and the Wilcoxon signed-rank $z$ of (arm minus baseline).

\section{Results on the real space}

\textbf{Nothing meaningfully beats random.} In absolute terms this is a space where random
search is essentially unbeatable. Across every budget, random and all four crossovers finish
within a few tenths of a percent of one another (Table~\ref{tab:scaled}), a gap comparable to
the benchmark's own training noise: the same architecture trained three times varies by a
similar amount. No crossover produces a gain a practitioner would notice. What the paired,
many-seed design buys is not a large effect but resolution: it lets us read a consistent
ordering inside that noise band, which is what the rest of this section reports.

\textbf{Structure does not help; at large budget it slips below random.}
Table~\ref{tab:scaled} scales to population $50$, budget $20{,}000$, $4000$ paired seeds. All
differences are small (within the noise floor above) but consistent across the $4000$ seeds.
Flat is best at every budget and beats each structured operator head-to-head; its small edge
over random shrinks as random saturates ($z: 48, 46, 29, 13$ across the four budgets) but stays
consistent ($2268/1322$ paired wins even at $20{,}000$). The structured operators worsen with
compute: axis crosses from $z=+30$ to $z=-28$ (losing to random $1230/2559$); path barely moves
across a fortyfold budget rise ($93.904$ to $93.913$) and ends at $z=-26$; aligned decays from
$z=42$ to $z=1.9$, to break-even. So the structure-blind operator wins throughout, the two most
structured (axis, path) fall below random at large budget, and the intermediate one ties it.
This is not a validity artifact: axis, aligned, and path in fact yield more valid offspring
($95$ to $97$ percent) than flat ($90$ to $96$ percent).

\begin{table}[t]
\centering
\caption{NAS-Bench-101, population $50$, elite $12$, $4000$ paired seeds. Mean best
test-accuracy (\%) and Wilcoxon signed-rank $z$ against random search (positive means above
random). All gaps are small (within the benchmark's training noise) but consistent across
$4000$ seeds. Flat is best at every budget; axis and path slip just below random at large
budget; aligned's edge decays to zero.}
\label{tab:scaled}
\begin{tabular}{lccccc}
\toprule
budget & random & flat ($z$) & axis ($z$) & aligned ($z$) & path ($z$) \\
\midrule
$500$   & $93.645$ & $\mathbf{93.998}\;(+48.0)$ & $93.806\;(+30.3)$ & $93.917\;(+41.8)$ & $93.904\;(+44.4)$ \\
$2000$  & $93.746$ & $\mathbf{94.072}\;(+45.8)$ & $93.872\;(+22.5)$ & $94.009\;(+39.8)$ & $93.911\;(+29.9)$ \\
$8000$  & $93.941$ & $\mathbf{94.085}\;(+29.2)$ & $93.883\;(-10.2)$ & $94.036\;(+19.7)$ & $93.912\;(-6.1)$ \\
$20000$ & $94.042$ & $\mathbf{94.088}\;(+12.5)$ & $93.886\;(-27.6)$ & $94.044\;(+1.9)$ & $93.913\;(-26.0)$ \\
\bottomrule
\end{tabular}
\end{table}

\textbf{Mechanism: productive diversity.} Two independent measurements locate the cause. First,
the final population's genotypic spread (mean pairwise distance over $21$ edges and $5$ ops)
shows two failure modes. The operators that end below random, axis and path, collapse the spread
($1.8$ and $1.3$ against flat's $2.1$ at $20{,}000$): they converge early, matching their
plateaus. But aligned keeps the most spread of all ($3.1$, above flat) and still only ties
random, because its variety never concentrates on good cells. Flat holds an intermediate,
useful spread and wins. (We had expected the cleaner story that diversity ranks the operators;
the measurement does not support it, and we report that.) Second, varying the population size at
fixed budget sweeps the diversity available to crossover, and the comparison moves with it
(Table~\ref{tab:pop}): at population $8$ every arm, flat included, loses to random ($z<0$); the
gain rises steadily until at population $64$ all four beat it. Crossover here beats random
exactly when it has enough diversity and an operator that turns diversity into good offspring
rather than collapsing or wasting it.

\begin{table}[t]
\centering
\caption{Population sweep at fixed budget $2000$, $600$ paired seeds (elite $=$ pop$/4$; the
random control is identical throughout, mean $93.743$). Mean best test-accuracy (\%) and
Wilcoxon $z$ vs random: at pop $8$ every arm loses to random; the gain rises with population.}
\label{tab:pop}
\begin{tabular}{lcccc}
\toprule
pop & flat ($z$) & axis ($z$) & aligned ($z$) & path ($z$) \\
\midrule
$8$  & $93.673\;(-4.5)$ & $93.652\;(-5.7)$ & $93.673\;(-4.2)$ & $93.526\;(-12.6)$ \\
$16$ & $93.791\;(+3.5)$ & $93.715\;(-1.8)$ & $93.756\;(+1.1)$ & $93.706\;(-2.7)$  \\
$32$ & $93.981\;(+14.6)$& $93.823\;(+5.9)$ & $93.898\;(+10.5)$& $93.840\;(+7.3)$  \\
$64$ & $\mathbf{94.122}\;(+19.3)$& $93.917\;(+11.6)$& $94.064\;(+17.2)$& $93.957\;(+14.2)$ \\
\bottomrule
\end{tabular}
\end{table}

\textbf{Why random is hard to beat here: good cells are common.} The small effect sizes have a
simple structural cause, visible in the benchmark's accuracy distribution. Good cells are
common: of the $423{,}624$ cells, $2.45\%$ ($10{,}378$) score within one point of the best
($94.32\%$), $17.7\%$ within two points, and the median cell already reaches about $90\%$. So a
single random draw is within one point of the best with probability $0.0245$, and the best of
just $100$ random draws is within one point of the optimum $92\%$ of the time. Random search
reaches the good region almost at once, and at that resolution there is nothing left for a
smarter search to win. The last sliver is the opposite case: only $0.09\%$ ($400$ cells) lie
within half a point of the best, so that final fraction of a percent is genuinely rare, and the
differences there are inside the benchmark's own training noise, so no method can resolve them
reliably. Between the good region, reached at once by sampling, and the true optimum, rare and
below the noise, there is no band where a directed search wins a real, repeatable margin. That
is why random search is not beaten here, and why crossover's edge is real as a direction but
very small in size.

\section{Relation to prior work}

Nothing here is new in framing. Building-block recombination and crossover are Holland
\cite{holland1975}; deceptive traps as their canonical test are Deb and Goldberg \cite{deb1993};
random search as a strong NAS baseline is \cite{bergstra2012,li2019}; mutation-only evolutionary
NAS is \cite{real2019}; the competing-conventions problem the aligned operator targets motivates
NEAT's historical markings \cite{neat2002} and He et al.'s modular inheritable crossover
\cite{modxover2021}; NAS-Bench-101 \cite{nasbench101} makes a study at this seed count possible.
The contribution is a clean, reproducible pairing: a positive control that shows exactly when
crossover's building-block advantage appears, and a real NAS space where it does not. Four
structure-respecting crossovers all lose (consistently, though by small margins) to plain
uniform crossover, two of them to random search, with a diversity mechanism that explains why.

\section{Discussion}

One rule ties the two testbeds together: crossover helps to the extent the space is built from
separable, recombinable blocks. The trap control makes the effect plain: crossover's edge over
mutation grows with the number of blocks until mutation cannot solve the problem. The real NAS
cell does not supply such structure along its natural dimensions. Recombining its wiring and
labeling, aligning nodes by role, and transplanting whole paths all lose to a structure-blind
uniform crossover, whose own edge over random comes only from diversity, which the structural
priors collapse or waste.

\textbf{On effect size.} None of the NAS-Bench-101 differences is large. Every method lands
within a few tenths of a percent of random, inside the benchmark's training noise, so no
crossover meaningfully beats random search. Our claims are about the consistent direction of
small effects resolved by thousands of paired seeds, not about their size. The practical message
is the familiar one, that random search is very hard to beat on this space, made sharper by the
fact that even the direction does not favor added structure.

\textbf{Limitations.} This is one tabulated cell space with a lookup fitness (no training-time
confound, but also no guarantee of transfer to larger or hierarchical spaces). The aligned
operator is one alignment, and a better one may exist; the burden is now on showing one that
beats uniform crossover, which four natural attempts do not. The natural next step is a larger,
hierarchical space where recombination has more room to matter.

\section*{Reproducibility}

Every number regenerates from \cite{repo} with one command. The trap control is
\texttt{bbtest.c} (\texttt{make bbtest}); the fitness table is extracted by
\texttt{nb101\_extract.c} (\texttt{make nb101\_extract}); the search, oracle, and all operators
are \texttt{validation/nb101.c} (\texttt{make nb101}), configurable through the environment
(\texttt{POP}, \texttt{ELITE}, \texttt{SEEDS}, \texttt{BUDGETS}) and self-testing its oracle
with \texttt{--selftest}. All build clean under \texttt{-std=c99 -pedantic} and run clean under
AddressSanitizer and UBSan; a fixed xorshift generator makes every run reproducible.

\begin{thebibliography}{99}
\small
\bibitem{thesis1997} V.~Gavrilov. \emph{Backpropagation Feed-Forward Neural Networks}.
  Seminar in Artificial Intelligence, Open University of Israel, 1997 (revised digital
  edition, Zenodo 2026). \url{https://doi.org/10.5281/zenodo.20450526}
\bibitem{holland1975} J.~H. Holland. \emph{Adaptation in Natural and Artificial Systems}.
  University of Michigan Press, 1975.
\bibitem{deb1993} K.~Deb, D.~E. Goldberg. Analyzing Deception in Trap Functions.
  \emph{Foundations of Genetic Algorithms} 2, 1993.
\bibitem{bergstra2012} J.~Bergstra, Y.~Bengio. Random Search for Hyper-Parameter
  Optimization. \emph{JMLR} 13, 2012.
\bibitem{zoph2017} B.~Zoph, Q.~V. Le. Neural Architecture Search with Reinforcement
  Learning. \emph{ICLR} 2017.
\bibitem{darts2019} H.~Liu, K.~Simonyan, Y.~Yang. DARTS: Differentiable Architecture
  Search. \emph{ICLR} 2019.
\bibitem{real2019} E.~Real, A.~Aggarwal, Y.~Huang, Q.~V. Le. Regularized Evolution for
  Image Classifier Architecture Search. \emph{AAAI} 2019.
\bibitem{elsken2019} T.~Elsken, J.~H. Metzen, F.~Hutter. Neural Architecture Search: A
  Survey. \emph{JMLR} 20, 2019.
\bibitem{li2019} L.~Li, A.~Talwalkar. Random Search and Reproducibility for Neural
  Architecture Search. \emph{UAI} 2019.
\bibitem{neat2002} K.~O. Stanley, R.~Miikkulainen. Evolving Neural Networks through
  Augmenting Topologies. \emph{Evolutionary Computation} 10(2), 2002.
\bibitem{modxover2021} C.~He, H.~Tan, S.~Huang, R.~Cheng. Efficient evolutionary neural
  architecture search by modular inheritable crossover. \emph{Swarm and Evolutionary
  Computation} 64, 2021.
\bibitem{nasbench101} C.~Ying, A.~Klein, E.~Real, E.~Christiansen, K.~Murphy, F.~Hutter.
  NAS-Bench-101: Towards Reproducible Neural Architecture Search. \emph{ICML} 2019.
\bibitem{repo} SMBPANN reference implementation.
  \url{https://github.com/Anode1/SMBPANN}.
\end{thebibliography}

\end{document}
